% Gabarit pandoc des corrigés publiés (brut et guidé). Contenu sans marque ni nom de site.
% Compilation : lualatex (latexmk -lualatex).
% Le Markdown reste la source unique : ne pas éditer ce fichier .tex à la main.
\documentclass[11pt,a4paper]{article}

% --- Moteur, polices (Latin Modern) et langue ---------------------------------
\usepackage{iftex}
\ifLuaTeX\else\errmessage{Ce document se compile avec lualatex}\fi
\usepackage{amsmath,amssymb,mathtools}
\usepackage{fontspec}
% mathrm=sym : \mathrm prend ses glyphes dans la police mathématique (ℓ, μ, Ω), comme KaTeX.
\usepackage[math-style=TeX,bold-style=TeX,mathrm=sym]{unicode-math}
% Repli pour les glyphes absents de Latin Modern (indices et exposants Unicode,
% flèches, filets des schémas en texte, symboles) : DejaVu, puis Latin Modern Math
% (ℓ, lettres mathématiques Unicode des légendes).
\directlua{
  luaotfload.add_fallback("repliTexte", {"DejaVuSerif:mode=harf;", "DejaVuSans:mode=harf;", "file:latinmodern-math.otf:mode=harf;"})
  luaotfload.add_fallback("repliMono", {"DejaVuSansMono:mode=harf;", "DejaVuSans:mode=harf;", "file:latinmodern-math.otf:mode=harf;"})
}
\setmainfont{Latin Modern Roman}[RawFeature={fallback=repliTexte}]
\setsansfont{Latin Modern Sans}[RawFeature={fallback=repliTexte}]
\setmonofont{Latin Modern Mono}[RawFeature={fallback=repliMono}]
\setmathfont{Latin Modern Math}
\usepackage[french]{babel}
\frenchsetup{SmallCapsFigTabCaptions=false}

% --- Mise en page --------------------------------------------------------------
\usepackage[a4paper,margin=2.2cm,headheight=14pt,footskip=1.2cm]{geometry}
\usepackage{microtype}
\setlength{\emergencystretch}{4em}
\tolerance=1500
\setlength{\parindent}{0pt}
\setlength{\parskip}{0.55\baselineskip plus 2pt minus 1pt}
\usepackage{xcolor}
% Gris nommés : babel-french rend « ! » actif dans le document.
\definecolor{gris2}{gray}{0.98}
\definecolor{gris3}{gray}{0.97}
\definecolor{gris5}{gray}{0.95}
\definecolor{gris35}{gray}{0.65}
\definecolor{gris40}{gray}{0.60}
\definecolor{gris45}{gray}{0.55}
\definecolor{gris50}{gray}{0.50}
\definecolor{gris55}{gray}{0.45}
\definecolor{gris60}{gray}{0.40}
\definecolor{gris65}{gray}{0.35}
\definecolor{gris70}{gray}{0.30}
\definecolor{gris80}{gray}{0.20}
\definecolor{gris85}{gray}{0.15}
\usepackage{graphicx}
\usepackage{float}
\usepackage{caption}
\captionsetup{font=small,labelformat=empty,justification=centering}
\usepackage{booktabs,longtable,array,multirow,calc}
\usepackage{siunitx}
\sisetup{locale=FR}
\usepackage[version=4]{mhchem}
\usepackage{listings}
\usepackage{fancyhdr}
\usepackage{lastpage}
\usepackage[most]{tcolorbox}
\usepackage{enumitem}
\setlist{nosep,leftmargin=*,topsep=2pt}
\usepackage[hidelinks,unicode]{hyperref}
\usepackage{bookmark}
\hypersetup{pdftitle={Centrale-Supélec MP-MPI-PC-PSI — Anglais 2026},pdfauthor={Rédigé par IA},pdfcreator={LaTeX},pdfsubject={Corrigé brut}}
\usepackage{xurl}

% --- Pandoc : listes, tableaux, figures ----------------------------------------
\providecommand{\passthrough}[1]{#1}
\providecommand{\tightlist}{\setlength{\itemsep}{0pt}\setlength{\parskip}{0pt}}

\makeatletter
\def\maxwidth{\ifdim\Gin@nat@width>\linewidth\linewidth\else\Gin@nat@width\fi}
\def\maxheight{\ifdim\Gin@nat@height>0.55\textheight 0.55\textheight\else\Gin@nat@height\fi}
\makeatother
\setkeys{Gin}{width=\maxwidth,height=\maxheight,keepaspectratio}
\setcounter{secnumdepth}{-\maxdimen}
\providecommand{\pandocbounded}[1]{#1}

% --- Code (listings) -----------------------------------------------------------
\lstset{
  basicstyle=\ttfamily\small,
  columns=fullflexible,
  keepspaces=true,
  breaklines=true,
  breakatwhitespace=false,
  postbreak=\mbox{\textcolor{gris50}{$\hookrightarrow$}\space},
  frame=single,
  rulecolor=\color{gris35},
  framesep=4pt,
  xleftmargin=5pt,xrightmargin=5pt,
  backgroundcolor=\color{gris3},
  keywordstyle=\bfseries,
  commentstyle=\color{gris65},
  stringstyle=\color{gris80},
  showstringspaces=false,
  tabsize=4,
  extendedchars=true,
  aboveskip=6pt,belowskip=6pt,
}

% --- Macros des corrigés (KaTeX ↔ LaTeX) ---------------------------------------
% Les corrigés sont écrits pour KaTeX : on définit ici ce que KaTeX connaît et
% que LaTeX n'a pas, sans jamais retoucher le Markdown.
\providecommand{\R}{\mathbb{R}}
\providecommand{\N}{\mathbb{N}}
\providecommand{\Z}{\mathbb{Z}}
\providecommand{\Q}{\mathbb{Q}}
\providecommand{\C}{\mathbb{C}}
\providecommand{\dd}{\mathrm{d}}
\providecommand{\llbracket}{\lBrack}
\providecommand{\rrbracket}{\rBrack}
\providecommand{\lBrace}{\lbrace}
\providecommand{\rBrace}{\rbrace}
\providecommand{\Box}{\square}
\providecommand{\Diamond}{\lozenge}
\providecommand{\varnothing}{\emptyset}
\providecommand{\argmax}{\operatorname*{arg\,max}}
\providecommand{\argmin}{\operatorname*{arg\,min}}
\renewcommand{\leq}{\leqslant}
\renewcommand{\geq}{\geqslant}

% Équation hors texte : réduite à la largeur de la ligne si elle déborde.
\newsavebox{\boiteequation}
\newcommand{\equationajustee}[1]{%
  \sbox{\boiteequation}{$\displaystyle #1$}%
  \ifdim\wd\boiteequation>\linewidth
    \[\resizebox{\linewidth}{!}{\usebox{\boiteequation}}\]%
  \else
    \[#1\]%
  \fi}

% --- Encadrés pédagogiques (sobres, lisibles en noir et blanc) -----------------
\tcbset{
  encadre/.style={
    enhanced,breakable,sharp corners,
    colback=white,colframe=gris55,
    boxrule=0.4pt,leftrule=2.2pt,
    left=6pt,right=6pt,top=3pt,bottom=3pt,
    before skip=6pt,after skip=6pt,
    fonttitle=\bfseries\small,coltitle=black,
    attach title to upper,after title={\quad},
    parbox=false,
  }
}
\newtcolorbox{indication}[1]{encadre,colframe=gris45,colback=gris2,title={Indication~#1}}
\newtcolorbox{methode}{encadre,colframe=gris70,title={Méthode}}
\newtcolorbox{piege}{encadre,colframe=gris85,boxrule=0.8pt,leftrule=2.2pt,title={Piège}}
\newtcolorbox{jury}{encadre,colframe=gris60,colback=gris5,title={Jury}}
\newtcolorbox{encadrenote}{encadre,colframe=gris40}

% --- Titres de section ---------------------------------------------------------
\usepackage{titlesec}
\titleformat{\section}{\Large\bfseries}{}{0pt}{}
\titleformat{\subsection}{\large\bfseries}{}{0pt}{}
\titleformat{\subsubsection}{\normalsize\bfseries}{}{0pt}{}
\titlespacing*{\section}{0pt}{14pt plus 4pt minus 2pt}{6pt}
\titlespacing*{\subsection}{0pt}{10pt plus 3pt minus 2pt}{4pt}
\titlespacing*{\subsubsection}{0pt}{8pt plus 2pt minus 2pt}{2pt}

% --- En-tête et pied de page ---------------------------------------------------
\pagestyle{fancy}
\fancyhf{}
\fancyhead[L]{\small\textsc{Corrigé brut} · Centrale-Supélec MP-MPI-PC-PSI}
\fancyhead[R]{\small Anglais 2026}
\fancyfoot[L]{\footnotesize Corrigé rédigé par IA, non officiel}
\fancyfoot[R]{\footnotesize \thepage\,/\,\pageref*{LastPage}}
\renewcommand{\headrulewidth}{0.4pt}
\renewcommand{\footrulewidth}{0pt}
\fancypagestyle{plain}{\fancyhf{}\fancyfoot[R]{\footnotesize \thepage\,/\,\pageref*{LastPage}}\fancyfoot[L]{\footnotesize Corrigé rédigé par IA, non officiel}\renewcommand{\headrulewidth}{0pt}}

\begin{document}
\thispagestyle{plain}
\begin{center}
  {\small\scshape Corrigé brut}\par\smallskip
  {\LARGE\bfseries Centrale-Supélec MP-MPI-PC-PSI — Anglais 2026\par}\medskip
  {\large Centrale-Supélec · MP-MPI-PC-PSI · Anglais · 2026\par}
\end{center}
\vspace{2pt}\hrule\vspace{8pt}

\emph{Corrigé rédigé par IA, non officiel. Signalez toute erreur.}

Épreuve de 4 heures, sans calculatrice ni système électronique. Consigne : rédiger en anglais une synthèse de 500 mots des quatre documents, avec un titre obligatoire ; indiquer à la fin le nombre de mots, titre inclus ; un écart de 10 \% est accepté, soit \textbf{450 à 550 mots}.

Documents : photographie de David Dawson (2001), Lucian Freud peignant le portrait d'Elizabeth II ; David Smith, \emph{The Guardian}, 31 décembre 2024 (Kim Sajet et la National Portrait Gallery de la Smithsonian) ; Kelly Grovier, \emph{BBC.com}, 1er août 2025 (la statue de la Liberté et le tableau d'Amy Sherald) ; Roger Kimball, \emph{The Spectator}, 13 août 2025 (le décret de Donald Trump sur la Smithsonian).

Convention de décompte retenue : un mot est une suite de caractères entre deux espaces ; un mot à trait d'union (\emph{ever-changing}), un nombre (\emph{1886}, \emph{250th}) et une forme à apostrophe (\emph{Smithsonian's}) comptent chacun pour un mot. Le titre est compté.

\textbf{Problématique} : portraits et symboles nationaux peuvent-ils encore dire la vérité d'une nation lorsque le pouvoir politique prétend en fixer le sens ?

\hypertarget{synthuxe8se}{%
\section{Synthèse}\label{synthuxe8se}}

\textbf{Painting the Nation: Truth, Power and the Portrait}

Who decides how a nation sees itself? A 2001 photograph by David Dawson shows Lucian Freud painting Queen Elizabeth II, a portrait now held by London's National Portrait Gallery. Three articles, from The Guardian, the BBC and The Spectator, move the question to Washington, where the Smithsonian's museums have become a battlefield during Donald Trump's second term. They ask whether portraits and national symbols can still tell the truth once politics claims them.

The documents first agree that a portrait is a construction rather than a mirror. For Kim Sajet, director of the Smithsonian's National Portrait Gallery, ``all portraiture is a lie'': every image results from a negotiation between the sitter, the artist, the patron and an ever-changing audience. Dawson's photograph makes this visible. The Queen sits smiling and composed, wearing a diadem, while the painter, seen from behind, studies her before a small canvas on which her face looks heavier and sterner. National symbols are no different. Kelly Grovier recalls that the Statue of Liberty drew on Roman, Greek and Egyptian sources and has divided Americans since 1886, when suffragettes found it ironic and conservatives feared it would attract migrants. Amy Sherald's painting of Liberty as a black transgender woman reopens that old fault line.

Hence, the institutions that display such images have become political targets. Sajet's gallery strives to remain ``above the fray'': its labels try ``not to editorialise'' and rest on historical fact, as the carefully worded caption of Trump's photograph shows. Roger Kimball sees things differently. Applauding Trump's March executive order, ``Restoring Truth and Sanity to American History'', he accuses the Smithsonian of spreading a ``divisive, race-centered ideology'' and welcomes an official review of its exhibition texts and curatorial procedures, meant to ``celebrate American exceptionalism'' before the country's 250th anniversary. Grovier shows what such pressure means in practice, in the very museum Sajet presented as neutral: fearing for their funding, the organisers of Sherald's exhibition had second thoughts about her painting, and the artist cancelled the whole show, citing ``a culture of censorship''. Both camps thus claim to defend the truth, but they disagree on who should define it.

Finally, the documents question whether meaning can be controlled at all. Kimball urges conservatives to ``commandeer'' the institutions of elite culture, as the left has done since the 1960s. Grovier replies that symbols escape such control: excluding Sherald's work has only increased its impact, for nothing is ``more visible than something hidden''. Sajet trusts time and distance. Her gallery presents every president from George Washington to Barack Obama and Donald Trump, so that visitors can keep the present in perspective, and Trump's official portrait, though already completed, will stay in storage until 2029. As a Nigerian-born Dutchwoman, she also values an outsider's ability to stand back from America.

In the end, these documents suggest that no portrait tells the whole truth, yet, as Picasso put it, art may be ``a lie that makes us realise truth''. The real danger lies less in this partiality than in any attempt to impose a single, official image of a nation.

\textbf{Number of words (title included): 518}

\end{document}
